\documentclass[12pt]{article}
\usepackage[utf8]{inputenc}
\usepackage[T1]{fontenc}
\usepackage{hyperref}
\title{Waluwarra Language Resources for K-12 Curriculum Development}
\author{Compiled by Claude Sonnet 5.5 (claude-sonnet-5-5) for Andr\'as Kornai}
\date{}
\begin{document}
\maketitle

\section{Introduction}
Waluwarra (also Warluwarra, Waluwara; ISO 639-3 wrb; Glottocode warl1256) is a Pama--Nyungan language of western Queensland, placed in the Southern Ngarna subgroup. It is extinct. Sources disagree on when, with dates ranging from 1990 to after 2005. Gavan Breen documented it, with three dialect groups: Kapula, Parnkara and Warluwara proper. Care is needed because the Wikipedia article ``Walubarra'' is about a different, Yidiny-speaking group near Cairns.

\section{Primary Educational Resources}
None found. No primer, textbook or school programme appeared in search results.

\section{Linguistic and Reference Materials}
\begin{itemize}
\item Wikipedia, ``Warluwarra language'': \url{https://en.wikipedia.org/wiki/Warluwarra_language}
\item AIATSIS Austlang entry G10 (the search summary linked Warluwarra to Breen's 1970 Monash MA thesis and the wordlist ASEDA 0253; the page itself was not opened): \url{https://collection.aiatsis.gov.au/austlang/language/g10}
\item AIATSIS Austlang entries for the related varieties Parnkarra (G41) and Kapula (G43): \url{https://collection.aiatsis.gov.au/austlang/language/g41} and \url{https://collection.aiatsis.gov.au/austlang/language/g43}
\item ASJP wordlist for Warluwara (Carrizo and Sidwell): \url{https://asjp.clld.org/languages/WARLUWARA}
\item Wikipedia, ``Ngarna languages'' (classification): \url{https://en.wikipedia.org/wiki/Ngarna_languages}
\end{itemize}

\section{Religious and Community Resources}
None found.

\section{Digital and Audio Resources}
The Endangered Languages Project has pages that may relate to this language, but the search results did not show their content: \url{https://endangeredlanguages.com/lang/3355/samples/10194}

\section{Curriculum Development Context}
There are no speakers, and only a short wordlist and Breen's thesis-level description appear to exist. A revitalization curriculum would have to be built from archival material at AIATSIS and with descendant communities, who should decide how it is used. Such a curriculum would be modest in scope.

\end{document}